\chapter{Two-Dimensional Superconformal Algebras}
\label{ch:two-dimensional-superconformal-algebras}

\ConventionsEuclideanCFT

\section{Spin and Sector Data}

The superconformal algebras in two dimensions are chiral symmetry algebras.
They are not, by themselves, definitions of a continuum quantum field theory.
The underlying QFT must supply local operator domains, correlation functions,
positivity when the theory is unitary, and sewing data.  This chapter isolates
the chiral algebraic layer because it is shared infrastructure for
two-dimensional CFT.  The dynamical construction of supersymmetric field
theories, including Landau--Ginzburg and gauged linear sigma models, belongs
to the supersymmetric-QFT development in Volume~VII.  The statements below are
therefore algebraic consequences and consistency tests for any such model,
not a replacement for the supersymmetric RG analysis.

\begin{definition}[Spin two-dimensional CFT datum]
\label{def:spin-two-dimensional-cft-datum}
A spin two-dimensional CFT datum consists of the following data.
\begin{enumerate}
\item An oriented two-dimensional Euclidean CFT \(\mathcal C\), with
holomorphic stress tensor \(T(z)\), antiholomorphic stress tensor
\(\bar T(\bar z)\), and central charges \(c,\bar c\).
\item For each spin surface \((\Sigma,\mathfrak s)\), correlation functions
of local fields which may be even or odd under the fermion-parity operator
\((-1)^F\).
\item Chiral state spaces on the circle decomposed into Neveu--Schwarz and
Ramond sectors,
\[
  \Hilb^{\mathrm{ch}}
  =
  \Hilb_{\mathrm{NS}}^{\mathrm{ch}}
  \oplus
  \Hilb_{\mathrm R}^{\mathrm{ch}},
\]
where a chiral fermionic field has antiperiodic boundary condition in the
NS sector and periodic boundary condition in the R sector.
\item A locality rule: even fields have ordinary single-valued OPEs with all
local fields, while odd chiral fields are mutually local only after the spin
structure and fermion-parity signs have been specified.
\end{enumerate}
The same chiral algebra may occur in several full spin CFTs, because the full
theory also requires an antiholomorphic sector, a choice of spin-structure
sum or projection, and modular/sewing data.
\end{definition}

Let \(a(z)\) be a chiral field of holomorphic weight \(h_a\).  In a sector
where \(a\) has monodromy \(e^{-2\pi \ii \nu_a}\) around the spatial circle,
we write
\begin{equation}
  a(z)=\sum_{r\in\mathbb Z+\nu_a} a_r z^{-r-h_a}.
  \label{eq:spin-cft-chiral-mode-expansion}
\end{equation}
For the supercurrent of an \(\mathcal N=1\) or \(\mathcal N=2\)
superconformal algebra, \(h_a=3/2\), \(\nu_a=1/2\) in the NS sector, and
\(\nu_a=0\) in the R sector.  The Ramond sector therefore contains
supercharge zero modes.  This is not a choice of notation; it is the operator
form of the difference between antiperiodic and periodic spin structures on
the cylinder.

\section{\texorpdfstring{\(\mathcal N=1\)}{N=1} Superconformal Algebra}

\begin{definition}[\(\mathcal N=1\) chiral superconformal datum]
\label{def:n1-superconformal-datum}
An \(\mathcal N=1\) chiral superconformal datum of central charge \(c\)
consists of a chiral stress tensor \(T(z)\), an odd chiral field \(G(z)\) of
weight \(3/2\), and a sector choice
\[
  r\in\mathbb Z+\frac12\quad(\mathrm{NS}),\qquad
  r\in\mathbb Z\quad(\mathrm R)
\]
for the modes
\[
  T(z)=\sum_{n\in\mathbb Z}L_n z^{-n-2},
  \qquad
  G(z)=\sum_{r}G_r z^{-r-3/2}.
\]
The singular OPEs are
\begin{align}
  T(z)T(w)
  &\sim
  \frac{c/2}{(z-w)^4}
  +\frac{2T(w)}{(z-w)^2}
  +\frac{\partial T(w)}{z-w},
  \label{eq:n1-tt-ope}\\
  T(z)G(w)
  &\sim
  \frac{(3/2)G(w)}{(z-w)^2}
  +\frac{\partial G(w)}{z-w},
  \label{eq:n1-tg-ope}\\
  G(z)G(w)
  &\sim
  \frac{2c/3}{(z-w)^3}
  +\frac{2T(w)}{z-w}.
  \label{eq:n1-gg-ope}
\end{align}
\end{definition}

\begin{proposition}[Mode algebra from the \(\mathcal N=1\) OPEs]
\label{prop:n1-mode-algebra}
The modes in Definition~\ref{def:n1-superconformal-datum} obey
\begin{align}
  [L_m,L_n]
  &=
  (m-n)L_{m+n}
  +\frac{c}{12}(m^3-m)\delta_{m+n,0},
  \label{eq:n1-virasoro-mode-algebra}\\
  [L_m,G_r]
  &=
  \left(\frac m2-r\right)G_{m+r},
  \label{eq:n1-lg-mode-algebra}\\
  \{G_r,G_s\}
  &=
  2L_{r+s}
  +\frac c3\left(r^2-\frac14\right)\delta_{r+s,0}.
  \label{eq:n1-gg-mode-algebra}
\end{align}
\end{proposition}

\begin{proof}
The Virasoro bracket is the stress-tensor calculation of
Proposition~\ref{prop:tt-ope-virasoro-algebra}.  For the second formula,
write
\[
  [L_m,G_r]
  =
  \oint_0\frac{\dd w}{2\pi\ii}w^{r+1/2}
  \oint_w\frac{\dd z}{2\pi\ii}z^{m+1}T(z)G(w).
\]
Using \eqref{eq:n1-tg-ope}, the inner residue is
\[
  (m+1)w^m\frac32G(w)+w^{m+1}\partial G(w).
\]
The outer residue of the first term is
\(\frac32(m+1)G_{m+r}\).  For the second term, integration by parts gives
\[
  \oint_0\frac{\dd w}{2\pi\ii}w^{m+r+3/2}\partial G(w)
  =
  -\left(m+r+\frac32\right)G_{m+r}.
\]
The coefficient is therefore \(m/2-r\).

For the anticommutator, the odd fields require the supercommutator contour.
The singular part \eqref{eq:n1-gg-ope} gives
\[
  \{G_r,G_s\}
  =
  \oint_0\frac{\dd w}{2\pi\ii}w^{s+1/2}
  \operatorname{Res}_{z=w}
  z^{r+1/2}
  \left(
    \frac{2c/3}{(z-w)^3}+\frac{2T(w)}{z-w}
  \right).
\]
The \(T\)-term gives \(2L_{r+s}\).  The cubic pole gives
\[
  \frac{2c}{3}\cdot
  \frac12
  \left(r+\frac12\right)\left(r-\frac12\right)
  \delta_{r+s,0}
  =
  \frac c3\left(r^2-\frac14\right)\delta_{r+s,0}.
\]
\end{proof}

\begin{corollary}[Ramond zero-mode bound]
\label{cor:n1-ramond-zero-mode-bound}
In a unitary Ramond representation of the \(\mathcal N=1\) algebra,
\[
  L_0-\frac{c}{24}\ge0.
  \label{eq:n1-ramond-bound}
\]
A Ramond ground state is a state annihilated by \(G_0\), equivalently a state
with \(L_0=c/24\).
\end{corollary}

\begin{proof}
Set \(r=s=0\) in \eqref{eq:n1-gg-mode-algebra}.  Then
\[
  \{G_0,G_0\}=2L_0-\frac c{12}
  =
  2\left(L_0-\frac c{24}\right).
\]
If the representation is unitary, \(G_0^\dagger=G_0\), so the expectation
value of \(\{G_0,G_0\}=2G_0^2\) on any vector is nonnegative.  Equality
holds precisely on the kernel of \(G_0\).
\end{proof}

\section{\texorpdfstring{\(\mathcal N=2\)}{N=2} Superconformal Algebra}

The \(\mathcal N=2\) algebra adds a holomorphic \(U(1)\) current.  This extra
current is the reason that spectral flow, chiral-primary bounds, and
topological twists are algebraic operations rather than model-dependent
accidents.

\begin{definition}[\(\mathcal N=2\) chiral superconformal datum]
\label{def:n2-superconformal-datum}
An \(\mathcal N=2\) chiral superconformal datum of central charge \(c\)
consists of a chiral stress tensor \(T(z)\), a weight-one even current
\(J(z)\), and two odd weight-\(3/2\) chiral fields \(G^+(z)\) and \(G^-(z)\).
Their modes are
\[
  J(z)=\sum_{n\in\mathbb Z}J_n z^{-n-1},
  \qquad
  G^\pm(z)=\sum_r G^\pm_r z^{-r-3/2},
\]
with \(r\in\mathbb Z+1/2\) in the NS sector and \(r\in\mathbb Z\) in the R
sector.  The singular OPEs, in addition to the Virasoro OPE for \(T\), are
\begin{align}
  T(z)J(w)
  &\sim
  \frac{J(w)}{(z-w)^2}
  +\frac{\partial J(w)}{z-w},
  \label{eq:n2-tj-ope}\\
  T(z)G^\pm(w)
  &\sim
  \frac{(3/2)G^\pm(w)}{(z-w)^2}
  +\frac{\partial G^\pm(w)}{z-w},
  \label{eq:n2-tg-ope}\\
  J(z)J(w)
  &\sim
  \frac{c/3}{(z-w)^2},
  \label{eq:n2-jj-ope}\\
  J(z)G^\pm(w)
  &\sim
  \pm\frac{G^\pm(w)}{z-w},
  \label{eq:n2-jg-ope}\\
  G^+(z)G^-(w)
  &\sim
  \frac{2c/3}{(z-w)^3}
  +\frac{2J(w)}{(z-w)^2}
  +\frac{2T(w)+\partial J(w)}{z-w},
  \label{eq:n2-gplus-gminus-ope}\\
  G^\pm(z)G^\pm(w)&\sim0.
  \label{eq:n2-same-sign-g-ope}
\end{align}
\end{definition}

\begin{proposition}[Mode algebra for \(\mathcal N=2\)]
\label{prop:n2-mode-algebra}
The modes in Definition~\ref{def:n2-superconformal-datum} obey
\begin{align}
  [L_m,J_n]&=-nJ_{m+n},
  \label{eq:n2-lj-mode-algebra}\\
  [J_m,J_n]&=\frac c3 m\delta_{m+n,0},
  \label{eq:n2-jj-mode-algebra}\\
  [L_m,G^\pm_r]&=\left(\frac m2-r\right)G^\pm_{m+r},
  \label{eq:n2-lg-mode-algebra}\\
  [J_m,G^\pm_r]&=\pm G^\pm_{m+r},
  \label{eq:n2-jg-mode-algebra}\\
  \{G^+_r,G^-_s\}
  &=
  2L_{r+s}
  +(r-s)J_{r+s}
  +\frac c3\left(r^2-\frac14\right)\delta_{r+s,0},
  \label{eq:n2-gplus-gminus-mode-algebra}\\
  \{G^\pm_r,G^\pm_s\}&=0.
  \label{eq:n2-same-sign-g-mode-algebra}
\end{align}
\end{proposition}

\begin{proof}
The formulas involving \(L_m\) are the primary-field residue calculation
used in Proposition~\ref{prop:n1-mode-algebra}, with weights \(1\) and
\(3/2\).  The current commutator follows from the double-pole residue in
\eqref{eq:n2-jj-ope}:
\[
  [J_m,J_n]
  =
  \oint_0\frac{\dd w}{2\pi\ii}w^n
  \operatorname{Res}_{z=w}
  z^m\frac{c/3}{(z-w)^2}
  =
  \frac c3 m\delta_{m+n,0}.
\]
Equation~\eqref{eq:n2-jg-mode-algebra} is the simple-pole residue of
\eqref{eq:n2-jg-ope}.  For the mixed supercurrent anticommutator, the
triple pole in \eqref{eq:n2-gplus-gminus-ope} gives
\(\frac c3(r^2-1/4)\delta_{r+s,0}\), exactly as in the
\(\mathcal N=1\) calculation.  The simple pole gives \(2L_{r+s}\).  The
double pole and the derivative term combine as
\[
  2\left(r+\frac12\right)J_{r+s}
  -
  (r+s+1)J_{r+s}
  =
  (r-s)J_{r+s}.
\]
\end{proof}

\begin{proposition}[Spectral flow automorphism]
\label{prop:n2-spectral-flow}
For any \(\eta\in\frac12\mathbb Z\), the assignment
\begin{align}
  L_n^{(\eta)}
  &=
  L_n+\eta J_n+\frac c6\eta^2\delta_{n,0},
  \label{eq:n2-spectral-flow-l}\\
  J_n^{(\eta)}
  &=
  J_n+\frac c3\eta\delta_{n,0},
  \label{eq:n2-spectral-flow-j}\\
  (G^\pm_r)^{(\eta)}
  &=
  G^\pm_{r\pm\eta}
  \label{eq:n2-spectral-flow-g}
\end{align}
preserves the \(\mathcal N=2\) mode algebra, with the same central charge
\(c\).  Integer \(\eta\) preserves NS and R sectors separately; half-integer
\(\eta\) exchanges them.
\end{proposition}

\begin{proof}
It is enough to check the brackets in
Proposition~\ref{prop:n2-mode-algebra}.  For example,
\[
  [J_m^{(\eta)},J_n^{(\eta)}]
  =
  [J_m,J_n]
  =
  \frac c3m\delta_{m+n,0},
\]
which is the same formula in the flowed generators because the central shifts
commute.  Also
\[
  [J_m^{(\eta)},(G^\pm_r)^{(\eta)}]
  =
  [J_m,G^\pm_{r\pm\eta}]
  =
  \pm G^\pm_{m+r\pm\eta}
  =
  \pm(G^\pm_{m+r})^{(\eta)}.
\]
For the mixed supercurrent relation,
\[
  \{G^+_{r+\eta},G^-_{s-\eta}\}
  =
  2L_{r+s}
  +(r-s+2\eta)J_{r+s}
  +\frac c3\left((r+\eta)^2-\frac14\right)\delta_{r+s,0}.
\]
Substituting
\[
  L_{r+s}=L^{(\eta)}_{r+s}
  -\eta J_{r+s}
  -\frac c6\eta^2\delta_{r+s,0},
  \qquad
  J_{r+s}=J^{(\eta)}_{r+s}
  -\frac c3\eta\delta_{r+s,0},
\]
the \(J\)-coefficient becomes \(r-s\), and the central coefficient reduces
to \(\frac c3(r^2-1/4)\).  The remaining brackets are checked in the same
residue algebra.
\end{proof}

If a simultaneous \(L_0,J_0\) eigenstate has eigenvalues \(h,q\), then
spectral flow sends them to
\begin{equation}
  h^{(\eta)}
  =
  h+\eta q+\frac c6\eta^2,
  \qquad
  q^{(\eta)}
  =
  q+\frac c3\eta .
  \label{eq:n2-spectral-flow-hq}
\end{equation}
This formula is often the cleanest way to see why NS chiral primaries and
Ramond ground states are the same representation-theoretic object viewed in
different spin sectors.

\begin{proposition}[NS unitarity bound and chiral primaries]
\label{prop:n2-chiral-primary-bound}
Let \(v\) be a simultaneous \(L_0,J_0\) eigenvector in a unitary NS
representation, with
\[
  L_0v=hv,\qquad J_0v=qv.
\]
Then
\[
  h\ge\frac{|q|}{2}.
  \label{eq:n2-ns-unitarity-bound}
\]
If \(q\ge0\), equality \(h=q/2\) holds precisely when
\[
  G^+_{-1/2}v=0.
  \label{eq:n2-chiral-primary-condition}
\]
If \(q\le0\), equality \(h=-q/2\) holds precisely when
\[
  G^-_{-1/2}v=0.
  \label{eq:n2-antichiral-primary-condition}
\]
\end{proposition}

\begin{proof}
Unitarity gives \((G^\pm_{-1/2})^\dagger=G^\mp_{1/2}\).  Taking expectation
values of \eqref{eq:n2-gplus-gminus-mode-algebra} gives
\[
  \|G^+_{-1/2}v\|^2
  =
  \langle v,\{G^-_{1/2},G^+_{-1/2}\}v\rangle
  =
  (2h-q)\|v\|^2,
\]
and
\[
  \|G^-_{-1/2}v\|^2
  =
  \langle v,\{G^+_{1/2},G^-_{-1/2}\}v\rangle
  =
  (2h+q)\|v\|^2.
\]
Both norms are nonnegative, so \(2h\ge q\) and \(2h\ge -q\).  Equality in
the corresponding inequality is equivalent to the vanishing of the
corresponding norm.
\end{proof}

\begin{corollary}[NS--R relation for chiral primaries]
\label{cor:n2-chiral-primary-ramond-ground-state}
If \(v\) is an NS chiral primary with \(h=q/2\), then spectral flow by
\(\eta=-1/2\) maps it to a Ramond state with
\[
  h^{(-1/2)}=\frac c{24}.
  \label{eq:n2-chiral-to-ramond-ground}
\]
Thus it saturates the Ramond ground-state energy.
\end{corollary}

\begin{proof}
Insert \(h=q/2\) and \(\eta=-1/2\) into
\eqref{eq:n2-spectral-flow-hq}:
\[
  h^{(-1/2)}
  =
  \frac q2-\frac q2+\frac c{24}
  =
  \frac c{24}.
\]
\end{proof}

\section{Protected Algebraic Tests for Supersymmetric Dynamics}

The full construction of two-dimensional \(\mathcal N=(2,2)\)
Landau--Ginzburg and gauged linear sigma-model flows belongs to the
supersymmetric-QFT chapters.  The part needed in the CFT volume is narrower:
the chiral-ring and central-charge data that any proposed
\(\mathcal N=2\) superconformal fixed point must reproduce.  We state this
as an algebraic interface between the CFT and supersymmetric-QFT volumes, not
as a proof of an infrared fixed point.

\begin{definition}[Quasihomogeneous Landau--Ginzburg chiral datum]
\label{def:quasihomogeneous-lg-chiral-datum}
Let \(X_1,\ldots,X_N\) be complex variables with rational weights
\(q_i\in(0,1/2]\).  A quasihomogeneous Landau--Ginzburg chiral datum consists
of a polynomial
\[
  W\in\mathbb C[X_1,\ldots,X_N]
\]
such that
\[
  W(\lambda^{q_1}X_1,\ldots,\lambda^{q_N}X_N)=\lambda W(X_1,\ldots,X_N)
  \qquad
  \lambda\in\mathbb C^\times,
  \label{eq:lg-quasihomogeneous-condition}
\]
and such that the critical locus
\[
  \partial_1W=\cdots=\partial_NW=0
\]
is isolated at the origin.  Its Jacobi ring is
\[
  \operatorname{Jac}(W)
  =
  \mathbb C[X_1,\ldots,X_N]/
  (\partial_1W,\ldots,\partial_NW).
  \label{eq:lg-jacobi-ring}
\]
\end{definition}

The \(U(1)\) charge of a monomial
\[
  X^a=X_1^{a_1}\cdots X_N^{a_N}
\]
is
\[
  Q_R(X^a)=\sum_{i=1}^N a_iq_i.
  \label{eq:lg-monomial-r-charge}
\]
If the datum is realized by a unitary compact \(\mathcal N=2\)
superconformal fixed point with these fields as chiral primaries, the
chiral-primary relation gives
\[
  h(X^a)=\frac12 Q_R(X^a).
  \label{eq:lg-chiral-weight-from-charge}
\]
The central charge predicted by the \(U(1)_R\) anomaly is
\begin{equation}
  c_{\mathrm{LG}}
  =
  3\sum_{i=1}^N (1-2q_i).
  \label{eq:lg-central-charge}
\end{equation}
In this chapter \eqref{eq:lg-central-charge} is used only as a protected
chiral-algebra consistency check.  The assertion that a given ultraviolet
Landau--Ginzburg functional flows to a particular continuum SCFT requires
the regulator and RG construction requested in the Volume~VII
supersymmetric-QFT program.

\begin{example}[\(A\)-series algebra]
\label{ex:a-series-lg-algebra}
For
\[
  W=X^{k+2},
  \qquad k\in\mathbb Z_{\ge0},
\]
the quasihomogeneous weight is \(q=1/(k+2)\).  The Jacobi ring is
\[
  \operatorname{Jac}(W)
  =
  \mathbb C[X]/(X^{k+1}),
  \label{eq:a-series-jacobi-ring}
\]
with basis \(1,X,\ldots,X^k\).  The central charge predicted by
\eqref{eq:lg-central-charge} is
\[
  c=\frac{3k}{k+2},
  \label{eq:a-series-lg-central-charge}
\]
and the basis element \(X^\ell\), \(0\le\ell\le k\), has
\[
  q_\ell=\frac{\ell}{k+2},
  \qquad
  h_\ell=\frac{\ell}{2(k+2)}
  \label{eq:a-series-lg-chiral-weights}
\]
when represented by a chiral primary.
\end{example}

\begin{example}[Quintic algebraic check]
\label{ex:quintic-lg-algebraic-check}
For
\[
  W=X_1^5+\cdots+X_5^5
\]
with all weights \(q_i=1/5\), the anomaly formula gives
\[
  c=3\cdot 5\left(1-\frac25\right)=9.
  \label{eq:quintic-lg-central-charge}
\]
This is an algebraic \(U(1)_R\)-anomaly calculation.  Interpreting it as a
Calabi--Yau sigma-model phase is not part of the present chapter.
\end{example}

\section{Small \texorpdfstring{\(\mathcal N=4\)}{N=4} Boundary}

For completeness, we record the boundary between what is used in this
monograph and the full small \(\mathcal N=4\) representation theory.  The
small \(\mathcal N=4\) algebra contains an affine \(\mathfrak{su}(2)\) current
algebra and four supercurrents.  It includes \(\mathcal N=2\) subalgebras
after choosing a Cartan \(U(1)\subset SU(2)_R\), but the representation
theory involves the full \(SU(2)_R\) spin.

\begin{definition}[Small \(\mathcal N=4\) chiral datum]
\label{def:small-n4-chiral-datum}
A small \(\mathcal N=4\) chiral datum at level \(k\in\mathbb Z_{\ge0}\)
consists of:
\begin{enumerate}
\item a stress tensor \(T(z)\) of central charge \(c=6k\);
\item affine \(\mathfrak{su}(2)\) currents \(J^a(z)\), \(a=1,2,3\), with
level \(k\);
\item four weight-\(3/2\) supercurrents transforming as two doublets under
\(\mathfrak{su}(2)\);
\item OPEs such that the anticommutator of opposite supercurrents closes on
\(T\), \(J^a\), and the central term, while the current OPEs act on the
supercurrents by the doublet representation.
\end{enumerate}
\end{definition}

Definition~\ref{def:small-n4-chiral-datum} records only the part of the
small-\(\mathcal N=4\) algebra used in this volume.  The present monograph
uses two consequences: the central charge is tied to the affine level by
\(c=6k\), and every Cartan choice supplies an \(\mathcal N=2\) subalgebra
whose spectral-flow operation is inherited from the affine \(SU(2)_R\)
structure.  A later chapter that needs small-\(\mathcal N=4\) multiplet
classification must state the full OPE tensor and derive the corresponding
unitarity bounds.

\begin{openproblem}[Intrinsic two-dimensional supersymmetric RG construction]
\label{op:intrinsic-2d-supersymmetric-rg-construction}
Coordinate the CFT-volume algebraic layer with the Volume~VII construction of
two-dimensional \(\mathcal N=(2,2)\) Landau--Ginzburg models, gauged linear
sigma models, and their proposed infrared SCFTs as intrinsic two-dimensional
QFTs with a regulator, supersymmetric Ward identities, renormalized
\(U(1)_R\) current, and a topology in which the infrared chiral ring and
central charge are obtained.  The algebraic checks in
Definition~\ref{def:quasihomogeneous-lg-chiral-datum} and
Examples~\ref{ex:a-series-lg-algebra}--\ref{ex:quintic-lg-algebraic-check}
are necessary consistency tests, not a construction of the continuum fixed
point.
\end{openproblem}
